
\documentclass{article}
\usepackage{colortbl}
\usepackage{makecell}
\usepackage{multirow}
\usepackage{supertabular}

\begin{document}

\newcounter{utterance}

\centering \large Interaction Transcript for game `eqbench', experiment `eq\_bench', episode 69 with reflexion{-}llama{-}3.1{-}8b{-}instruct.
\vspace{24pt}

{ \footnotesize  \setcounter{utterance}{1}
\setlength{\tabcolsep}{0pt}
\begin{supertabular}{c@{$\;$}|p{.15\linewidth}@{}p{.15\linewidth}p{.15\linewidth}p{.15\linewidth}p{.15\linewidth}p{.15\linewidth}}
    \# & \multicolumn{2}{c}{Player} && \multicolumn{2}{c}{Game Master} \\
    \hline

    \theutterance \stepcounter{utterance}  
    & & & \multicolumn{4}{p{0.6\linewidth}}{
        \cellcolor[rgb]{0.9,0.9,0.9}{
            \makecell[{{p{\linewidth}}}]{
                \texttt{\tiny{[P1$\langle$GM]}}
                \texttt{You have attempted this same task before and did not succeed. These are the plans you wrote after those attempts.} \\
\texttt{Trial \#1: New plan: I overestimated Jacques' emotional response to Olivia's witty remark about his cooking. While it's clear that Jacques is defensive and insulted, I didn't consider the possibility that he might be more hurt by the personal attack on his ego. Next time, I will pay closer attention to the dialogue and consider the character's emotional vulnerabilities. I will also score the emotions more carefully, taking into account the context and the character's previous statements. For example, Jacques' initial jabs at Olivia suggest that he is already feeling confident and assertive, which might make him less likely to feel embarrassment.} \\
\\ 
\texttt{Trial \#2: New plan: I overestimated the intensity of Jacques' emotions, particularly embarrassment. While Jacques is clearly defensive and insulted, his reaction seems more driven by anger and a desire to assert his ego. Next time, I will pay closer attention to the dialogue and consider the character's emotional vulnerabilities, as well as the context and tone of their statements. I will also score the emotions more carefully, taking into account the character's previous statements and body language. For example, Jacques' initial jabs at Olivia suggest that he is already feeling confident and assertive, which might make him less likely to feel embarrassment. I will also consider the possibility that Jacques' anger might be a mask for other emotions, such as hurt or frustration. To improve my scoring, I will ask myself questions like: What are the character's emotional vulnerabilities? How do their previous statements and actions suggest they might react to certain situations? What is the tone and context of their dialogue? By considering these factors, I can develop a more nuanced understanding of the character's emotions and improve my scoring.} \\
\\ 
\\ 
\\ 
\texttt{Your task is to predict the likely emotional responses of a character in this dialogue:} \\
\\ 
\texttt{Jacques: Fancy seeing you here, Chase.} \\
\texttt{Olivia: The feeling is mutual, Duval. } \\
\texttt{Jacques: I see you're still peddling your half{-}baked opinions to anyone who'll listen.} \\
\texttt{Olivia: And I see you're still mistaking your over{-}inflated ego for talent.} \\
\texttt{Jacques: Talent? I've got more talent in my little finger than you have in your entire body.} \\
\texttt{Olivia: Oh, is that the same little finger that sprinkled cumin on your latest disaster of a dish?} \\
\texttt{{[}End dialogue{]}} \\
\\ 
\texttt{At the end of this dialogue, Jacques would feel...} \\
\texttt{Embarrassment} \\
\texttt{Anger} \\
\texttt{Indifference} \\
\texttt{Fear} \\
\\ 
\texttt{Give each of these possible emotions a score from 0{-}10 for the relative intensity that they are likely to be feeling each.} \\
\\ 
\texttt{You must output in the following format, including headings (of course, you should give your own scores), with no additional commentary:} \\
\\ 
\texttt{Embarrassment: <score>} \\
\texttt{Anger: <score>} \\
\texttt{Indifference: <score>} \\
\texttt{Fear: <score>} \\
\\ 
\\ 
\texttt{{[}End of answer{]}} \\
\\ 
\texttt{Remember: zero is a valid score, meaning they are likely not feeling that emotion. You must score at least one emotion > 0.} \\
\\ 
\texttt{Your answer:} \\
            }
        }
    }
     \\ \\

    \theutterance \stepcounter{utterance}  
    & \multicolumn{4}{p{0.6\linewidth}}{
        \cellcolor[rgb]{0.561,0.737,0.561}{
            \makecell[{{p{\linewidth}}}]{
                \texttt{\tiny{[P1$\rangle$GM]}}
                \texttt{Embarrassment: 2} \\
\texttt{Anger: 9} \\
\texttt{Indifference: 0} \\
\texttt{Fear: 0} \\
            }
        }
    }
    & & \\ \\

    \theutterance \stepcounter{utterance}  
    & & & \multicolumn{2}{p{0.3\linewidth}}{
        \cellcolor[rgb]{0.9,0.9,0.9}{
            \makecell[{{p{\linewidth}}}]{
                \texttt{\tiny{[GM$|$GM]}}
                \texttt{Embarrassment: 2} \\
\texttt{Anger: 9} \\
\texttt{Indifference: 0} \\
\texttt{Fear: 0} \\
            }
        }
    }
    & & \\ \\

    \theutterance \stepcounter{utterance}  
    & & & \multicolumn{2}{p{0.3\linewidth}}{
        \cellcolor[rgb]{0.9,0.9,0.9}{
            \makecell[{{p{\linewidth}}}]{
                \texttt{\tiny{[GM$|$GM]}}
                \texttt{Embarrassment: 7} \\
\texttt{Anger: 8} \\
\texttt{Indifference: 0} \\
\texttt{Fear: 0} \\
            }
        }
    }
    & & \\ \\

    \theutterance \stepcounter{utterance}  
    & & & \multicolumn{2}{p{0.3\linewidth}}{
        \cellcolor[rgb]{0.9,0.9,0.9}{
            \makecell[{{p{\linewidth}}}]{
                \texttt{\tiny{[GM$|$GM]}}
                \texttt{game\_result = LOSE} \\
            }
        }
    }
    & & \\ \\

    \theutterance \stepcounter{utterance}  
    & & \multicolumn{4}{p{0.6\linewidth}}{
        \cellcolor[rgb]{0.95,0.95,0.95}{
            \makecell[{{p{\linewidth}}}]{
                \texttt{\tiny{[P1$\leftrightarrow$P]}}
                \texttt{Post{-}episode Reflexion (not sent to GM):} \\
\texttt{New plan: I overestimated the intensity of Jacques' emotions, particularly embarrassment, and didn't consider the possibility that his anger might be a mask for other emotions. Next time, I will pay closer attention to the dialogue and consider the character's emotional vulnerabilities, as well as the context and tone of their statements. I will also score the emotions more carefully, taking into account the character's previous statements and body language. To improve my scoring, I will ask myself questions like: What are the character's emotional vulnerabilities? How do their previous statements and actions suggest they might react to certain situations? What is the tone and context of their dialogue? I will also consider the possibility that Jacques' anger might be a result of his hurt ego, and that his confidence and assertiveness might be a defense mechanism. By considering these factors, I can develop a more nuanced understanding of the character's emotions and improve my scoring.} \\
            }
        }
    }
    & \\ \\

\end{supertabular}
}

\end{document}
